\begin{textsample}{POS Dim 3 – ms – Score 10.00 – ms/ms\_inf\_e\_ef\_25.txt}  \label{ex:f3_pos_236}
7.2 Direitos de Aprendizagem As Diretrizes Curriculares Nacionais para a Educação \textbf{Infantil} – DCNEI definem como eixos norteadores das práticas pedagógicas as interações e a \textbf{brincadeira}.
Com o objetivo de garantir à \textbf{criança} acesso a processos de apropriação e articulação de conhecimentos e aprendizagens de diferentes linguagens, assim como o direito à proteção, saúde, liberdade, confiança, ao respeito, à dignidade, convivência e a interação com outras \textbf{crianças} e \textbf{adultos}, foram estabelecidos na Base Nacional Comum Curricular ( BNCC ) seis direitos de aprendizagem e desenvolvimento que deverão permear as vivências de todas as \textbf{crianças} brasileiras.
São eles : - Conviver com outras \textbf{crianças} e \textbf{adultos}, em \textbf{pequenos} e grandes grupos, utilizando diferentes linguagens, ampliando o conhecimento de si e do outro, o respeito em relação à cultura e às diferenças entre as pessoas. - \textbf{Brincar} cotidianamente de diversas formas em diferentes espaços e tempos, com diferentes parceiros ( \textbf{crianças} e \textbf{adultos} ), ampliando e diversificando seu acesso a produções culturais, seus conhecimentos, sua imaginação, sua criatividade, suas experiências emocionais, corporais, \textbf{sensoriais}, expressivas, cognitivas, sociais e relacionais. - Participar ativamente, com \textbf{adultos} e outras \textbf{crianças}, tanto do planejamento da gestão da escola e das atividades propostas pelo educador quanto da realização das atividades da vida cotidiana, tais como a escolha das \textbf{brincadeiras}, dos materiais e dos ambientes, desenvolvendo diferentes linguagens e elaborando conhecimentos, decidindo e se posicionando. - Explorar movimentos, \textbf{gestos}, \textbf{sons}, formas, \textbf{texturas}, cores, palavras, emoções, transformações, relacionamentos, histórias, objetos, elementos da natureza, na escola e fora dela, ampliando seus saberes sobre a cultura em suas diversas modalidades : as artes, a escrita, a ciência e a tecnologia. - Expressar, como sujeito dialógico, criativo e sensível, suas necessidades, emoções, sentimentos, dúvidas, hipóteses, descobertas, opiniões, questionamentos, por meio de diferentes linguagens. - Conhecer -se e construir sua identidade pessoal, social e cultural, constituindo uma imagem positiva de si e de seus grupos de pertencimento, nas diversas experiências de \textbf{cuidados}, interações, \textbf{brincadeiras} e linguagens vivenciadas na instituição escolar e em seu contexto familiar e comunitário.
Os Direitos de Aprendizagem e Desenvolvimento estão articulados aos três princípios elencados nas DCNEI, quais sejam, princípios éticos : da autonomia, responsabilidade, da solidariedade e do respeito ao bem comum, ao meio ambiente e às diferentes culturas, identidades e singularidades, correspondem aos direitos de conviver e conhecer -se ; políticos : dos direitos de cidadania, do exercício da criticidade e do respeito à ordem democrática, correspondentes aos direitos de expressar -se e de participar ; estéticos : da sensibilidade, criatividade, ludicidade e da liberdade de expressão nas diferentes manifestações artísticas e culturais, referentes aos direitos de \textbf{brincar} e explorar.
Portanto, esses princípios devem nortear todas as ações dos profissionais e as experiências vivenciadas pelas \textbf{crianças}.
Na garantia desses direitos é necessário considerar que as \textbf{crianças} são diferentes, possuem culturas diversas, trazendo com elas valores, atitudes, procedimentos e vivências de realidades diversificadas, portanto devem ser respeitados o tempo, o modo de aprender e de se expressar de cada uma delas.
A intencionalidade educativa na Educação \textbf{Infantil} consiste em ações planejadas, nas quais os direitos de aprendizagem devem ser assegurados, considerando as relações e interações das \textbf{crianças} com os saberes e práticas sociais e os conhecimentos sistematizados, sem perder de vista as competências gerais da Educação Básica.

% matched lemmas: adulto, brincadeira, brincar, criança, cuidado, gesto, infantil, pequeno, sensorial, som, textura
\end{textsample}
